% ch22_natural_transformations_2morphisms.tex — Volume II Chapter 10

\chapter{Natural Transformations as 2-Morphisms}

\begin{overviewbox}
Chapter 4 introduced natural transformations and showed that they organize
functors into a more structured whole. But we only used one operation on
natural transformations: vertical composition. There are actually three
canonical ways to compose natural transformations --- \term{vertical},
\term{horizontal}, and \term{whiskering} --- and they interact through a
single remarkable law: the \term{interchange law}.

The structure formed by categories, functors, and natural transformations,
with all three compositions, is called the \term{2-category} $\mathbf{Cat}$.
Understanding $\mathbf{Cat}$ as a 2-category is not just organizational:
it reveals that the three compositions are not ad hoc, but arise from the
fact that $\mathbf{Cat}$ has \emph{two levels of morphism}, each with its
own composition. Natural transformations are the \term{2-morphisms}.

The interchange law is the key coherence axiom: it says that composing in
the horizontal direction and then vertically gives the same result as
composing vertically and then horizontally. Without this, the two
compositions would conflict; with it, they mesh into a unified structure.

This chapter prepares the ground for enriched category theory (Volume III)
by making explicit the 2-categorical structure that has been implicit
throughout the book.
\end{overviewbox}


% ═══════════════════════════════════════════════════════════════════════════
\begin{nameorigin}
\term{2-category} was introduced by \emph{Jean Bénabou} (1963) and
developed by Kelly, Street, and many others. The terminology comes
from the dimensional count: a $2$-category has objects ($0$-cells),
$1$-morphisms ($1$-cells), and $2$-morphisms ($2$-cells), with each
``cell'' a higher-dimensional generalization of an arrow. The
generalization continues: $n$-categories have $k$-cells for $0 \leq k
\leq n$; $\infty$-categories (Vol IV) have cells at every level.

\term{Whiskering} (composing a natural transformation with a functor on
one side) is named for the visual picture: you imagine the natural
transformation as a ``hair'' or ``whisker'' attached to the side, then
``brushing it'' along the functor's input or output. The whisker
metaphor is Mac Lane's; the term is universal in 2-categorical
literature.

\term{Interchange law} expresses the compatibility of vertical and
horizontal composition of $2$-cells; it is the central coherence axiom
of a $2$-category. The terminology is from algebra (interchange of
operations) and dates to the early 1960s.
\end{nameorigin}

\section{The 2-Category \texorpdfstring{$\mathbf{Cat}$}{Cat}}
\label{sec:2cat-cat}
% ═══════════════════════════════════════════════════════════════════════════

\subsection{Objects, 1-Morphisms, 2-Morphisms}

The structure $\mathbf{Cat}$ has three levels:
\begin{itemize}
  \item \textbf{Objects (0-cells):} small categories $\mathcal{A},
        \mathcal{B}, \mathcal{C}, \ldots$
  \item \textbf{1-morphisms (1-cells):} functors $F, G : \mathcal{A} \to
        \mathcal{B}$.
  \item \textbf{2-morphisms (2-cells):} natural transformations
        $\alpha : F \Rightarrow G$ between parallel functors.
\end{itemize}

A \term{2-category} requires that between any two 0-cells $\mathcal{A}$
and $\mathcal{B}$, the collection of 1-cells forms a \emph{category}
--- a category $\mathbf{Cat}(\mathcal{A}, \mathcal{B})$ whose objects
are functors $\mathcal{A} \to \mathcal{B}$ and whose morphisms are
natural transformations. The composition in this hom-category is called
\term{vertical composition}.

Additionally, the 1-cell composition (functor composition) must interact
with the 2-cells (natural transformations) coherently. This interaction
is captured by \term{horizontal composition} and \term{whiskering}.

\begin{howtosay}
\begin{itemize}
  \item ``$\alpha : F \Rightarrow G$'' reads ``$\alpha$ is a 2-cell from
        $F$ to $G$,'' or ``$\alpha$ goes between $F$ and $G$.''
  \item We draw 2-cells as double arrows: $\Rightarrow$.
  \item The hom-category $\mathbf{Cat}(\mathcal{A}, \mathcal{B})$ is also
        written $[\mathcal{A}, \mathcal{B}]$ or $\mathcal{B}^\mathcal{A}$.
\end{itemize}
\end{howtosay}

\subsection{Globular Diagrams}

It helps to draw the three cell levels as nested diagrams:
\begin{itemize}
  \item 0-cells (categories) are points: $\bullet$.
  \item 1-cells (functors) are arrows between points: $\bullet \to \bullet$.
  \item 2-cells (natural transformations) are regions bounded by two
        parallel arrows:
\end{itemize}
\[
  \begin{tikzcd}
    \mathcal{A}
      \arrow[r, bend left=30, "F", ""{name=U}]
      \arrow[r, bend right=30, "G"', ""{name=D}]
    & \mathcal{B}
    \arrow[Rightarrow, from=U, to=D, "\alpha"]
  \end{tikzcd}
\]
This picture makes it visual that natural transformations ``fill in''
the region between two parallel functors.


% ═══════════════════════════════════════════════════════════════════════════
\section{Vertical Composition}
\label{sec:vertical}
% ═══════════════════════════════════════════════════════════════════════════

\subsection{Definition}

Given natural transformations $\alpha : F \Rightarrow G$ and $\beta :
G \Rightarrow H$ (both $\mathcal{A} \to \mathcal{B}$), their
\term{vertical composite} $\beta \cdot \alpha : F \Rightarrow H$ is
defined componentwise:
\[
  (\beta \cdot \alpha)_c = \beta_c \circ \alpha_c : Fc \to Gc \to Hc
  \qquad \text{for each } c \in \mathcal{A}.
\]

\begin{dialog}
Vertical composition stacks natural transformations on top of each other:
two consecutive transformations between the same pair of categories
collapse into one. The name ``vertical'' comes from the globular picture:
you are composing within a single column (between $\mathcal{A}$ and
$\mathcal{B}$), stacking transformation regions vertically.
\end{dialog}

The naturality of $\beta \cdot \alpha$ is immediate: for any $f: c \to c'$,
\[
  H(f) \circ (\beta \cdot \alpha)_c = H(f) \circ \beta_c \circ \alpha_c
  = \beta_{c'} \circ G(f) \circ \alpha_c = \beta_{c'} \circ \alpha_{c'} \circ F(f)
  = (\beta \cdot \alpha)_{c'} \circ F(f).
\]

The identity for vertical composition is the \term{identity natural
transformation} $\mathrm{id}_F : F \Rightarrow F$, with
$(\mathrm{id}_F)_c = \mathrm{id}_{Fc}$.

Vertical composition is associative (since composition in $\mathcal{B}$
is), so $\mathbf{Cat}(\mathcal{A}, \mathcal{B})$ is indeed a category.

\subsection{The Vertical Composition Diagram}

\[
  \begin{tikzcd}
    \mathcal{A}
      \arrow[r, bend left=50, "F", ""{name=F}]
      \arrow[r, "G"{description}, ""{name=G1, below}, ""{name=G2, above}]
      \arrow[r, bend right=50, "H"', ""{name=H}]
    & \mathcal{B}
    \arrow[Rightarrow, from=F, to=G2, "\alpha"]
    \arrow[Rightarrow, from=G1, to=H, "\beta"]
  \end{tikzcd}
  \quad = \quad
  \begin{tikzcd}
    \mathcal{A}
      \arrow[r, bend left=30, "F", ""{name=U}]
      \arrow[r, bend right=30, "H"', ""{name=D}]
    & \mathcal{B}
    \arrow[Rightarrow, from=U, to=D, "\beta \cdot \alpha"]
  \end{tikzcd}
\]


% ═══════════════════════════════════════════════════════════════════════════
\section{Horizontal Composition}
\label{sec:horizontal}
% ═══════════════════════════════════════════════════════════════════════════

\subsection{The Setup}

Vertical composition works within a single hom-category. Horizontal
composition works \emph{across} hom-categories. Suppose we have:
\[
  \begin{tikzcd}
    \mathcal{A}
      \arrow[r, bend left=25, "F", ""{name=F1}]
      \arrow[r, bend right=25, "G"', ""{name=G1}]
    & \mathcal{B}
      \arrow[r, bend left=25, "H", ""{name=H1}]
      \arrow[r, bend right=25, "K"', ""{name=K1}]
    & \mathcal{C}
    \arrow[Rightarrow, from=F1, to=G1, "\alpha"]
    \arrow[Rightarrow, from=H1, to=K1, "\beta"]
  \end{tikzcd}
\]
Here $\alpha: F \Rightarrow G: \mathcal{A} \to \mathcal{B}$ and
$\beta: H \Rightarrow K: \mathcal{B} \to \mathcal{C}$ are ``side by
side.'' Their horizontal composite is a natural transformation
$\beta * \alpha: H \circ F \Rightarrow K \circ G: \mathcal{A} \to
\mathcal{C}$.

\subsection{The Formula (Two Versions)}

At a component $a \in \mathcal{A}$, the horizontal composite
$(\beta * \alpha)_a : HFa \to KGa$ is defined by either of two equal
formulas:
\begin{align*}
  (\beta * \alpha)_a
    &= K(\alpha_a) \circ \beta_{Fa}
    & &\text{(apply $\beta$ first, then $K$ on $\alpha$)} \\
    &= \beta_{Ga} \circ H(\alpha_a)
    & &\text{(apply $H$ on $\alpha$ first, then $\beta$).}
\end{align*}

\begin{theorem}
The two formulas agree: $K(\alpha_a) \circ \beta_{Fa} = \beta_{Ga} \circ H(\alpha_a)$.
\end{theorem}

\begin{proof}
\begin{prooflog}
\proofstep{Consider the morphism $\alpha_a : Fa \to Ga$ in $\mathcal{B}$}{This is the $a$-component of $\alpha: F \Rightarrow G$}
\proofstep{$\beta: H \Rightarrow K$ is a natural transformation, so it is natural at the morphism $\alpha_a$}{Naturality of $\beta$ at $\alpha_a : Fa \to Ga$ gives: $K(\alpha_a) \circ \beta_{Fa} = \beta_{Ga} \circ H(\alpha_a)$}
\proofstep{This is exactly the required equation}{The two formulas for $(\beta * \alpha)_a$ are equal by the naturality square of $\beta$ at $\alpha_a$}
\end{prooflog}
\end{proof}

The naturality of $\beta * \alpha$ (as a transformation $HF \Rightarrow KG$)
follows from the naturalities of $\alpha$ and $\beta$ and functoriality of
$H, K$.

\subsection{Horizontal Composition is Functorial}

The assignment $(\alpha, \beta) \mapsto \beta * \alpha$ is the action on
morphisms of a functor
\[
  * : \mathbf{Cat}(\mathcal{A}, \mathcal{B}) \times
        \mathbf{Cat}(\mathcal{B}, \mathcal{C}) \;\to\;
        \mathbf{Cat}(\mathcal{A}, \mathcal{C}).
\]
This is the \emph{composition functor} for the 2-category $\mathbf{Cat}$.
Preservation of identity transformations:
$\mathrm{id}_H * \mathrm{id}_F = \mathrm{id}_{HF}$, since
$(\mathrm{id}_H * \mathrm{id}_F)_a = \mathrm{id}_{HFa}$.


% ═══════════════════════════════════════════════════════════════════════════
\section{Whiskering}
\label{sec:whiskering}
% ═══════════════════════════════════════════════════════════════════════════

\subsection{Whiskering as a Special Case of Horizontal Composition}

\term{Whiskering} is the horizontal composition of a natural transformation
with an identity natural transformation, or equivalently, applying a
functor to a natural transformation.

\textbf{Right-whiskering.}\quad Given $\alpha: F \Rightarrow G$ (functors
$\mathcal{A} \to \mathcal{B}$) and $H: \mathcal{C} \to \mathcal{A}$, define
\[
  \alpha H : F \circ H \Rightarrow G \circ H
  \qquad (\alpha H)_c = \alpha_{Hc} \;\; \text{for } c \in \mathcal{C}.
\]
Right-whiskering by $H$ precomposes with $H$.

\textbf{Left-whiskering.}\quad Given $\alpha: F \Rightarrow G$ (functors
$\mathcal{A} \to \mathcal{B}$) and $K: \mathcal{B} \to \mathcal{C}$, define
\[
  K\alpha : K \circ F \Rightarrow K \circ G
  \qquad (K\alpha)_a = K(\alpha_a) \;\; \text{for } a \in \mathcal{A}.
\]
Left-whiskering by $K$ postcomposes with $K$.

\begin{howtosay}
\begin{itemize}
  \item ``$\alpha H$'' reads ``$\alpha$ right-whiskered by $H$'' or
        ``$\alpha$ at $H$'' (the subscript $H$ reminds us it acts as a
        right argument).
  \item ``$K\alpha$'' reads ``$\alpha$ left-whiskered by $K$'' or
        ``$K$ applied to $\alpha$.''
  \item In index notation: $(K\alpha)_a = K(\alpha_a)$, which looks just
        like applying the functor $K$ to each component.
\end{itemize}
\end{howtosay}

\subsection{The Monad Axioms via Whiskering}

Whiskering gives an elegant way to state the monad axioms. For a monad
$(T, \eta, \mu)$ on $\mathcal{C}$:
\begin{align*}
  \mu \circ \mu T &= \mu \circ T\mu
    & &\text{(associativity: whiskering $\mu$ on both sides of $T$)} \\
  \mu \circ \eta T &= \mathrm{id}_T = \mu \circ T\eta
    & &\text{(unit laws: whiskering $\eta$ on both sides)}
\end{align*}
Here $\mu T = \mu \cdot T$ means right-whiskering $\mu$ by $T$, and
$T\mu$ means left-whiskering $\mu$ by $T$.

\subsection{Adjunction Unit and Counit via Whiskering}

The monad multiplication from an adjunction is $\mu = G\varepsilon F$:
left-whisker the counit $\varepsilon : FG \Rightarrow \mathrm{Id}$ by $G$
to get $G\varepsilon: GFG \Rightarrow G$, then right-whisker by $F$ to
get $G\varepsilon F : GFGF \Rightarrow GF$. The two triangle identities
for an adjunction are:
\[
  \varepsilon F \cdot F\eta = \mathrm{id}_F
  \qquad \text{and} \qquad
  G\varepsilon \cdot \eta G = \mathrm{id}_G
\]
--- both expressed using whiskering.


% ═══════════════════════════════════════════════════════════════════════════
\section{The Interchange Law}
\label{sec:interchange}
% ═══════════════════════════════════════════════════════════════════════════

\subsection{Statement}

The \term{interchange law} is the core compatibility axiom between vertical
and horizontal composition. Suppose we have four natural transformations
arranged in a $2 \times 2$ grid:
\[
  \begin{tikzcd}
    \mathcal{A}
      \arrow[r, bend left=40, "F", ""{name=F1}]
      \arrow[r, "G"{description}, ""{name=G1, below}, ""{name=G2, above}]
      \arrow[r, bend right=40, "H"', ""{name=H1}]
    & \mathcal{B}
      \arrow[r, bend left=40, "J", ""{name=J1}]
      \arrow[r, "K"{description}, ""{name=K1, below}, ""{name=K2, above}]
      \arrow[r, bend right=40, "L"', ""{name=L1}]
    & \mathcal{C}
    \arrow[Rightarrow, from=F1, to=G2, "\alpha"]
    \arrow[Rightarrow, from=G1, to=H1, "\beta"]
    \arrow[Rightarrow, from=J1, to=K2, "\gamma"]
    \arrow[Rightarrow, from=K1, to=L1, "\delta"]
  \end{tikzcd}
\]
where $\alpha: F \Rightarrow G$, $\beta: G \Rightarrow H$,
$\gamma: J \Rightarrow K$, $\delta: K \Rightarrow L$.

\begin{theorem}[Interchange Law]
\[
  (\delta \cdot \gamma) * (\beta \cdot \alpha)
  \;=\; (\delta * \beta) \cdot (\gamma * \alpha).
\]
Reading this: ``horizontal composite of vertical composites equals
vertical composite of horizontal composites.''
\end{theorem}

\begin{proof}
\begin{prooflog}
\proofstep{Evaluate both sides at a component $a \in \mathcal{A}$}{We need: $[(\delta \cdot \gamma) * (\beta \cdot \alpha)]_a = [(\delta * \beta) \cdot (\gamma * \alpha)]_a$}
\proofstep{Left side using horizontal formula}{$((\delta \cdot \gamma) * (\beta \cdot \alpha))_a = L(\beta \cdot \alpha)_a \circ (\delta \cdot \gamma)_{Fa} = L(\beta_a \circ \alpha_a) \circ \delta_{Fa} \circ \gamma_{Fa}$}
\proofstep{Right side: $(\gamma * \alpha)_a = K\alpha_a \circ \gamma_{Fa}$ and $(\delta * \beta)_a = L\beta_a \circ \delta_{Ga}$}{So $[(\delta * \beta) \cdot (\gamma * \alpha)]_a = (L\beta_a \circ \delta_{Ga}) \circ (K\alpha_a \circ \gamma_{Fa})$}
\proofstep{Rewrite right side: $L\beta_a \circ \delta_{Ga} \circ K\alpha_a \circ \gamma_{Fa}$}{Insert $\delta_{Ga} \circ K\alpha_a$; naturality of $\delta$ at $\alpha_a: Fa \to Ga$ gives $\delta_{Ga} \circ K\alpha_a = L\alpha_a \circ \delta_{Fa}$}
\proofstep{Right side becomes $L\beta_a \circ L\alpha_a \circ \delta_{Fa} \circ \gamma_{Fa}$}{Using $L\beta_a \circ L\alpha_a = L(\beta_a \circ \alpha_a)$ (functoriality of $L$)}
\proofstep{Right side is $L(\beta_a \circ \alpha_a) \circ \delta_{Fa} \circ \gamma_{Fa}$}{This equals the left side; done}
\end{prooflog}
\end{proof}

\subsection{Why the Interchange Law Matters}

The interchange law is not just a technical identity --- it is the
\emph{reason} that the three compositions are compatible. Without it, one
could not freely rewrite a 2-categorical expression by changing the order
of vertical and horizontal operations. The interchange law says:
\begin{quote}
\emph{It does not matter whether you compose vertically first and then
horizontally, or horizontally first and then vertically.}
\end{quote}
This is precisely the 2-dimensional analogue of the axiom that in a
monoid, the binary product is well-defined.

\begin{dialog}
Think of the $2 \times 2$ grid of natural transformations. You can
``collapse'' it into a single transformation either by:
\begin{enumerate}
  \item First composing each row horizontally (getting two transformations),
        then composing those vertically.
  \item First composing each column vertically (getting two transformations),
        then composing those horizontally.
\end{enumerate}
The interchange law says you get the same answer either way.
\end{dialog}


% ═══════════════════════════════════════════════════════════════════════════
\section{Strict 2-Categories and Bicategories}
\label{sec:bicategories}
% ═══════════════════════════════════════════════════════════════════════════

\subsection{Strict 2-Categories}

A \term{strict 2-category} is a category enriched over the category
$\mathbf{Cat}$ of small categories: it has a set of 0-cells, a hom-category
between each pair of 0-cells, and composition functors satisfying
associativity and unit laws \emph{strictly} (as equations, not just up to
isomorphism).

$\mathbf{Cat}$ itself (with small categories as objects) is a strict
2-category. The hom-category $\mathbf{Cat}(\mathcal{A}, \mathcal{B})$
is the functor category $[\mathcal{A}, \mathcal{B}]$; functor composition
is strictly associative; the interchange law holds as a strict equality.

Other examples of strict 2-categories:
\begin{itemize}
  \item $\mathbf{2Cat}$: strict 2-categories, strict 2-functors, strict
        natural transformations.
  \item $\mathbf{MonCat}$: monoidal categories, strong monoidal functors,
        monoidal natural transformations.
  \item $R\text{-}\mathbf{Bimod}$: rings as objects, bimodules as
        1-morphisms, bimodule maps as 2-morphisms.
\end{itemize}

\subsection{Bicategories: Weakening the Axioms}

In many naturally occurring examples, 1-cell composition is associative
only up to coherent isomorphism, not strictly. A \term{bicategory}
(or \term{weak 2-category}) relaxes this:
\begin{itemize}
  \item Composition of 1-cells: $c \circ (b \circ a) \cong (c \circ b)
        \circ a$ via an \term{associator} 2-cell $\alpha_{a,b,c}$.
  \item Unit laws: $a \circ \mathrm{id} \cong a$ and $\mathrm{id} \circ a
        \cong a$ via \term{unitor} 2-cells.
  \item These 2-cell isomorphisms must satisfy coherence conditions
        (the pentagon axiom for associators, the triangle axiom for unitors).
\end{itemize}

\begin{example}{Bimodules as a Bicategory}
Objects: rings. 1-morphisms: $(R,S)$-bimodules ($R\text{-}S$-bimodules as
morphisms from $R$ to $S$). 2-morphisms: bimodule maps. Composition:
tensor product $M \otimes_S N$. This is a bicategory: the tensor product
is associative only up to a canonical isomorphism, not strictly.
\end{example}

\begin{example}{Spans as a Bicategory}
Objects: sets. 1-morphisms $A \to B$: spans, i.e., diagrams $A \leftarrow
X \rightarrow B$. 2-morphisms: maps of spans. Composition: pullback. This
is a bicategory because pullback is only associative up to isomorphism.
\end{example}

A key theorem (Mac Lane's coherence theorem, generalized) says:
every bicategory is biequivalent to a strict 2-category. So, up to
biequivalence, the weak and strict notions coincide. But in practice,
the weak version is more natural for many examples.


% ═══════════════════════════════════════════════════════════════════════════
\section{Exercises}
\label{sec:2cat-exercises}
% ═══════════════════════════════════════════════════════════════════════════

\begin{exercisesbox}
\begin{qa}
\qaitem{What are the three levels of $\mathbf{Cat}$ as a 2-category?}{0-cells (objects): small categories; 1-cells: functors; 2-cells: natural transformations.}
\qaitem{Define vertical composition of natural transformations $\alpha: F \Rightarrow G$ and $\beta: G \Rightarrow H$.}{$(\beta \cdot \alpha)_c = \beta_c \circ \alpha_c$ for each $c$. It is the composition in the hom-category $\mathbf{Cat}(\mathcal{A}, \mathcal{B})$.}
\qaitem{Define horizontal composition of $\alpha: F \Rightarrow G : \mathcal{A} \to \mathcal{B}$ and $\beta: H \Rightarrow K : \mathcal{B} \to \mathcal{C}$.}{$(\beta * \alpha)_a = K(\alpha_a) \circ \beta_{Fa} = \beta_{Ga} \circ H(\alpha_a)$. Both formulas are equal by naturality of $\beta$.}
\qaitem{Why do the two formulas for horizontal composition agree?}{Because $\beta$ is natural at $\alpha_a: Fa \to Ga$: the naturality square for $\beta$ at this morphism gives exactly $K(\alpha_a) \circ \beta_{Fa} = \beta_{Ga} \circ H(\alpha_a)$.}
\qaitem{State the interchange law.}{$(\delta \cdot \gamma) * (\beta \cdot \alpha) = (\delta * \beta) \cdot (\gamma * \alpha)$: for a $2\times 2$ grid of composable natural transformations.}
\qaitem{What is right-whiskering $\alpha H$ for $\alpha: F \Rightarrow G$ and $H: \mathcal{C} \to \mathcal{A}$?}{$\alpha H : FH \Rightarrow GH$ with $(\alpha H)_c = \alpha_{Hc}$.}
\qaitem{What is left-whiskering $K\alpha$ for $\alpha: F \Rightarrow G$ and $K: \mathcal{B} \to \mathcal{C}$?}{$K\alpha : KF \Rightarrow KG$ with $(K\alpha)_a = K(\alpha_a)$.}
\qaitem{How are whiskering and horizontal composition related?}{Whiskering is horizontal composition with an identity transformation: $K\alpha = \mathrm{id}_K * \alpha$ and $\alpha H = \alpha * \mathrm{id}_H$.}
\qaitem{Express the monad associativity axiom $\mu \circ \mu T = \mu \circ T\mu$ using whiskering notation.}{$\mu \cdot (\mu \cdot T) = \mu \cdot (T \cdot \mu)$ where $\mu T$ is $\mu$ right-whiskered by $T$ (i.e., $(\mu T)_c = \mu_{Tc}$) and $T\mu$ is $T$ left-whiskered with $\mu$ (i.e., $(T\mu)_c = T(\mu_c)$).}
\qaitem{What are the triangle identities for $F \dashv G$ written using whiskering?}{$\varepsilon F \cdot F\eta = \mathrm{id}_F$ and $G\varepsilon \cdot \eta G = \mathrm{id}_G$.}
\qaitem{What is a strict 2-category?}{A category enriched over $\mathbf{Cat}$: 0-cells, hom-categories between pairs of 0-cells, and strictly associative and unital composition functors satisfying interchange.}
\qaitem{What is a bicategory?}{A weakened 2-category where 1-cell composition is associative only up to coherent associator 2-cells satisfying the pentagon axiom, and unital only up to coherent unitor 2-cells satisfying the triangle axiom.}
\qaitem{Give an example of a bicategory that is not a strict 2-category.}{Rings and bimodules: composition (tensor product) is only associative up to canonical isomorphism, not strictly.}
\qaitem{What does the identity natural transformation at $F$ look like?}{$(\mathrm{id}_F)_c = \mathrm{id}_{Fc}$: the identity morphism at $Fc$ for each $c$. It is the identity for vertical composition.}
\qaitem{What is the horizontal composite of two identity transformations $\mathrm{id}_H * \mathrm{id}_F$?}{$\mathrm{id}_{HF}$: the identity transformation at the composite functor $HF$.}
\qaitem{In a strict 2-category, how many independent axioms govern the interactions of 0-, 1-, and 2-cells?}{The 1-cells form a category (associativity, units). The 2-cells between each pair of 1-cells form a category (associativity, units for vertical composition). Horizontal composition is functorial. And the interchange law holds. These account for all the axioms.}
\qaitem{What does ``$\mathbf{Cat}$ enriched over $\mathbf{Cat}$'' mean?}{Each hom-set $\mathbf{Cat}(\mathcal{A}, \mathcal{B})$ is not just a set but a category (the functor category). The composition map is a functor between these categories, not just a function of sets.}
\qaitem{State Mac Lane's coherence theorem for bicategories.}{Every bicategory is biequivalent to a strict 2-category.}
\qaitem{What is the vertical identity for the hom-category $\mathbf{Cat}(\mathcal{A}, \mathcal{B})$?}{The identity natural transformation $\mathrm{id}_F$ at each functor $F$.}
\qaitem{For a 2-category, what axiom ensures that composing 2-cells horizontally and then vertically gives the same result as composing vertically and then horizontally?}{The interchange law: $(\delta \cdot \gamma) * (\beta \cdot \alpha) = (\delta * \beta) \cdot (\gamma * \alpha)$.}
\end{qa}
\end{exercisesbox}


% ═══════════════════════════════════════════════════════════════════════════
\begin{summarybox}
\textbf{Three compositions.}\quad Natural transformations support three
compositions: \emph{vertical} ($\beta \cdot \alpha$: stacking within one
hom-category), \emph{horizontal} ($\beta * \alpha$: placing side by side
across composable hom-categories), and \emph{whiskering} (a special case
of horizontal composition involving an identity).

\textbf{Horizontal composition formula.}\quad
$(\beta * \alpha)_a = K(\alpha_a) \circ \beta_{Fa} = \beta_{Ga} \circ H(\alpha_a)$:
the two formulas agree by naturality of $\beta$. This is a micro-proof of
the interchange law in disguise.

\textbf{The interchange law.}\quad
$(\delta \cdot \gamma) * (\beta \cdot \alpha) = (\delta * \beta) \cdot (\gamma * \alpha)$:
the law that reconciles horizontal and vertical composition. It is proved
using naturality at a component and functoriality.

\textbf{2-category \textbf{Cat}.}\quad The triple (categories, functors,
natural transformations) with these compositions forms a strict 2-category.
This is the prototype for all 2-categorical thinking.

\textbf{Bicategories.}\quad The weak version of a 2-category, where
1-cell composition is only associative up to coherent isomorphism. Examples:
bimodules (tensor product), spans (pullback). Every bicategory is
biequivalent to a strict 2-category.
\end{summarybox}


% ═══════════════════════════════════════════════════════════════════════════
\section{Formal Restatement}
\label{sec:2cat-formal}
% ═══════════════════════════════════════════════════════════════════════════

\begin{formalrestatement}

\textbf{2-Category.}\quad A strict 2-category $\mathcal{K}$ consists of:
a set of 0-cells; for each pair $(A,B)$ of 0-cells, a hom-category
$\mathcal{K}(A,B)$ (whose objects are 1-cells, morphisms are 2-cells, and
composition is vertical); for each triple $(A,B,C)$, a composition functor
$\circ: \mathcal{K}(B,C) \times \mathcal{K}(A,B) \to \mathcal{K}(A,C)$
that is strictly associative and unital.

\textbf{Vertical composition.}\quad The composition in the hom-category:
if $\alpha: f \Rightarrow g$ and $\beta: g \Rightarrow h$ in
$\mathcal{K}(A,B)$, then $\beta \cdot \alpha : f \Rightarrow h$.

\textbf{Horizontal composition.}\quad The action of the composition
functor on 2-cells: if $\alpha: f \Rightarrow g$ in $\mathcal{K}(A,B)$
and $\beta: h \Rightarrow k$ in $\mathcal{K}(B,C)$, then
$\beta * \alpha: hf \Rightarrow kg$ in $\mathcal{K}(A,C)$.

\textbf{Interchange law.}\quad Since $\circ$ is a functor:
$(\delta \cdot \gamma) * (\beta \cdot \alpha) = (\delta * \beta) \cdot
(\gamma * \alpha)$. This is functoriality of $\circ$ applied to the pair
of composable pairs $((\gamma, \alpha), (\delta, \beta))$.

\textbf{Whiskering.}\quad Special cases of horizontal composition:
right-whiskering by a 1-cell $f: A \to B$ gives $(\alpha f)_x =
\alpha_{fx}$; left-whiskering by $g: B \to C$ gives $(g\alpha)_x =
g(\alpha_x)$.

\textbf{$\mathbf{Cat}$ as a 2-category.}\quad Objects: small categories.
Hom-categories $\mathbf{Cat}(\mathcal{A}, \mathcal{B}) = [\mathcal{A},
\mathcal{B}]$: functor categories. Composition: functor composition (on
objects) and horizontal composition of natural transformations (on
morphisms in the hom-categories). Strictly associative.

\textbf{Bicategory.}\quad As above but with associator isomorphisms
$a_{f,g,h}: (fg)h \xRightarrow{\cong} f(gh)$ and unitor isomorphisms,
satisfying the pentagon and triangle coherence conditions.

\end{formalrestatement}
